% sec3_5.tex — Section 3.5: Large-Sample Properties of GMM

The GMM formula \eqref{eq:3.4.8} defines GMM estimators, which are a set of estimators,
each indexed by the weighting matrix $\hat{\mathbf{W}}$.  You will be delighted to know that
\emph{every} estimator to be introduced in the next few chapters is a GMM estimator for some
choice of $\hat{\mathbf{W}}$.  The task of this section is to develop the large-sample theory for the
GMM estimator for any given choice of $\hat{\mathbf{W}}$, which can be carried out quite easily
with the techniques of the previous chapter.  The first half of this section extends
Propositions~2.1--2.4 of Chapter~2 to GMM estimators.  One issue that did not arise
in those propositions is which GMM estimators should be preferred to other GMM
estimators.  This is a question of choosing $\hat{\mathbf{W}}$ optimally and will be taken up in the
latter part of this section.

%% ─── Asymptotic Distribution of the GMM Estimator ──────────────────────────
\subsection*{Asymptotic Distribution of the GMM Estimator}

The large-sample theory for $\hat{\boldsymbol{\delta}}(\hat{\mathbf{W}})$ valid for any given choice of
$\hat{\mathbf{W}}$ is

\begin{proposition}[asymptotic distribution of the GMM estimator]
\label{prop:3.1}
\begin{enumerate}[label=(\alph*)]
\item \emph{(Consistency)} Under Assumptions~3.1--3.4,
  $\plim_{n\to\infty} \hat{\boldsymbol{\delta}}(\hat{\mathbf{W}}) = \boldsymbol{\delta}$.

\item \emph{(Asymptotic Normality)} If Assumption~3.3 is strengthened as
  Assumption~3.5, then
  \[
    \sqrt{n}\,(\hat{\boldsymbol{\delta}}(\hat{\mathbf{W}}) - \boldsymbol{\delta})
    \dto N\bigl(\mathbf{0},\; \avar(\hat{\boldsymbol{\delta}}(\hat{\mathbf{W}}))\bigr)
    \quad \text{as } n \to \infty,
  \]
  where
  \begin{equation}
    \avar(\hat{\boldsymbol{\delta}}(\hat{\mathbf{W}}))
    = (\boldsymbol{\Sigma}_{\mathbf{xz}}' \mathbf{W} \boldsymbol{\Sigma}_{\mathbf{xz}})^{-1}\,
      \boldsymbol{\Sigma}_{\mathbf{xz}}' \mathbf{W} \mathbf{S} \mathbf{W}
      \boldsymbol{\Sigma}_{\mathbf{xz}}\,
      (\boldsymbol{\Sigma}_{\mathbf{xz}}' \mathbf{W} \boldsymbol{\Sigma}_{\mathbf{xz}})^{-1}.
    \label{eq:3.5.1}
  \end{equation}
  (Recall: $\boldsymbol{\Sigma}_{\mathbf{xz}} \equiv \E(\mathbf{x}_i \mathbf{z}_i')$,
  $\mathbf{S} = \E(\mathbf{g}_i \mathbf{g}_i') = \E(\varepsilon_i^2\, \mathbf{x}_i \mathbf{x}_i')$,
  $\mathbf{W} \equiv \plim\, \hat{\mathbf{W}}$.)

\item \emph{(Consistent Estimate of $\avar(\hat{\boldsymbol{\delta}}(\hat{\mathbf{W}}))$)}
  Suppose there is available a consistent estimator, $\hat{\mathbf{S}}$, of $\mathbf{S}$
  $(K \times K)$.  Then, under Assumption~3.2,
  $\avar(\hat{\boldsymbol{\delta}}(\hat{\mathbf{W}}))$ is consistently estimated by
  \begin{equation}
    \widehat{\avar(\hat{\boldsymbol{\delta}}(\hat{\mathbf{W}}))}
    \equiv (\mathbf{S}_{\mathbf{xz}}' \hat{\mathbf{W}} \mathbf{S}_{\mathbf{xz}})^{-1}\,
      \mathbf{S}_{\mathbf{xz}}' \hat{\mathbf{W}} \hat{\mathbf{S}} \hat{\mathbf{W}}
      \mathbf{S}_{\mathbf{xz}}\,
      (\mathbf{S}_{\mathbf{xz}}' \hat{\mathbf{W}} \mathbf{S}_{\mathbf{xz}})^{-1},
    \label{eq:3.5.2}
  \end{equation}
  where $\mathbf{S}_{\mathbf{xz}}$ is the sample mean of $\mathbf{x}_i \mathbf{z}_i'$:
  \[
    \underset{(K \times L)}{\mathbf{S}_{\mathbf{xz}}}
    \equiv \frac{1}{n} \sum_{i=1}^{n} \mathbf{x}_i \mathbf{z}_i'.
  \]
\end{enumerate}
\end{proposition}

The ugly looking expression for the asymptotic variance will become much prettier
when we choose the weighting matrix optimally.  If you have gone through the
proof of Proposition~2.1, you should find it a breeze to prove Proposition~3.1.  The
key observations are:

\begin{enumerate}
\item[(1)] $\mathbf{S}_{\mathbf{xz}} \pto \boldsymbol{\Sigma}_{\mathbf{xz}}$
  \quad (by ergodic stationarity)

\item[(2)] $\bar{\mathbf{g}}$
  $\bigl(\equiv \frac{1}{n} \sum_{i=1}^{n} \mathbf{g}_i\bigr) \pto \mathbf{0}$
  \quad (by ergodic stationarity and the orthogonality conditions)

\item[(3)] $\sqrt{n}\, \bar{\mathbf{g}} \dto N(\mathbf{0}, \mathbf{S})$
  \quad (by Assumption~3.5).
\end{enumerate}

Consistency immediately follows if you use (1), (2), and Lemma~2.3(a) on the
expression for the sampling error \eqref{eq:3.4.11}.  To prove asymptotic normality, multiply
both sides of \eqref{eq:3.4.11} by $\sqrt{n}$ to obtain
\begin{equation}
  \sqrt{n}\,(\hat{\boldsymbol{\delta}}(\hat{\mathbf{W}}) - \boldsymbol{\delta})
  = (\mathbf{S}_{\mathbf{xz}}' \hat{\mathbf{W}} \mathbf{S}_{\mathbf{xz}})^{-1}\,
    \mathbf{S}_{\mathbf{xz}}' \hat{\mathbf{W}}\, \sqrt{n}\, \bar{\mathbf{g}}
  \label{eq:3.5.3}
\end{equation}
and use (3), Lemma~2.3(a), and Lemma~2.4(c).  Part~(c) of Proposition~3.1 follows
immediately from Lemma~2.3(a).

%% ─── Estimation of Error Variance ──────────────────────────────────────────
\subsection*{Estimation of Error Variance}

In Section~2.3, we proved the consistency of the OLS estimator, $s^2$, of the error
variance.  As was noted there, the result holds as long as the residual is from some
consistent estimator of the coefficient vector.  The same is true here.

\begin{proposition}[consistent estimation of error variance]
\label{prop:3.2}
\emph{For any consistent estimator, $\hat{\boldsymbol{\delta}}$, of $\boldsymbol{\delta}$, define
$\hat{\varepsilon}_i \equiv y_i - \mathbf{z}_i' \hat{\boldsymbol{\delta}}$.  Under Assumptions~3.1, 3.2,
plus the assumption that $\E(\mathbf{z}_i \mathbf{z}_i')$ (second moments of the regressors) exists
and is finite,}
\[
  \frac{1}{n} \sum_{i=1}^{n} \hat{\varepsilon}_i^2 \pto \E(\varepsilon_i^2),
\]
\emph{provided $\E(\varepsilon_i^2)$ exists and is finite.}
\end{proposition}

The proof is very similar to the proof of Proposition~2.2.  The relationship between
$\hat{\varepsilon}_i$ and $\varepsilon_i$ is given by
\begin{equation}
  \hat{\varepsilon}_i \equiv y_i - \mathbf{z}_i' \hat{\boldsymbol{\delta}}
  = \varepsilon_i - \mathbf{z}_i'(\hat{\boldsymbol{\delta}} - \boldsymbol{\delta}),
  \label{eq:3.5.4}
\end{equation}
so that
\begin{equation}
  \hat{\varepsilon}_i^2 = \varepsilon_i^2
  - 2(\hat{\boldsymbol{\delta}} - \boldsymbol{\delta})' \mathbf{z}_i \cdot \varepsilon_i
  + (\hat{\boldsymbol{\delta}} - \boldsymbol{\delta})' \mathbf{z}_i \mathbf{z}_i'
    (\hat{\boldsymbol{\delta}} - \boldsymbol{\delta}).
  \label{eq:3.5.5}
\end{equation}
Summing over $i$, we obtain
\begin{equation}
  \frac{1}{n} \sum_{i=1}^{n} \hat{\varepsilon}_i^2
  = \frac{1}{n} \sum_{i=1}^{n} \varepsilon_i^2
  - 2(\hat{\boldsymbol{\delta}} - \boldsymbol{\delta})'
    \frac{1}{n} \sum_{i=1}^{n} \mathbf{z}_i \cdot \varepsilon_i
  + (\hat{\boldsymbol{\delta}} - \boldsymbol{\delta})'
    \Bigl(\frac{1}{n} \sum_{i=1}^{n} \mathbf{z}_i \mathbf{z}_i'\Bigr)
    (\hat{\boldsymbol{\delta}} - \boldsymbol{\delta}).
  \label{eq:3.5.6}
\end{equation}
As usual, $\frac{1}{n} \sum_i \varepsilon_i^2 \pto \E(\varepsilon_i^2)$.
Since $\hat{\boldsymbol{\delta}}$ is consistent and
$\frac{1}{n} \sum_i \mathbf{z}_i \mathbf{z}_i'$ converges in probability to some finite matrix by assumption,
the last term vanishes.  It is easy to show that $\E(\mathbf{z}_i \cdot \varepsilon_i)$ exists and is
finite.\footnote{By the Cauchy--Schwarz inequality,
$\E(|z_{i\ell}\,\varepsilon_i|) \leq \sqrt{\E(z_{i\ell}^2) \cdot \E(\varepsilon_i^2)}$,
where $z_{i\ell}$ is the $\ell$-th element of $\mathbf{z}_i$.  Both $\E(z_{i\ell}^2)$ and
$\E(\varepsilon_i^2)$ are finite by assumption.}  Then, by ergodic stationarity,
\[
  \frac{1}{n} \sum_i \mathbf{z}_i \cdot \varepsilon_i \pto \text{some finite vector.}
\]
So the second term on the RHS of \eqref{eq:3.5.6}, too, vanishes.

%% ─── Hypothesis Testing ────────────────────────────────────────────────────
\subsection*{Hypothesis Testing}

It should now be routine to derive from Proposition~3.1(b) and 3.1(c) the following.

\begin{proposition}[robust $t$-ratio and Wald statistics]
\label{prop:3.3}
\emph{Suppose Assumptions~3.1--3.5 hold, and suppose there is available a consistent estimate
$\hat{\mathbf{S}}$ of $\mathbf{S}$ $(\equiv \avar(\bar{\mathbf{g}}) = \E(\mathbf{g}_i \mathbf{g}_i'))$.
Let}
\[
  \widehat{\avar(\hat{\boldsymbol{\delta}}(\hat{\mathbf{W}}))}
  \equiv (\mathbf{S}_{\mathbf{xz}}' \hat{\mathbf{W}} \mathbf{S}_{\mathbf{xz}})^{-1}\,
    \mathbf{S}_{\mathbf{xz}}' \hat{\mathbf{W}} \hat{\mathbf{S}} \hat{\mathbf{W}}
    \mathbf{S}_{\mathbf{xz}}\,
    (\mathbf{S}_{\mathbf{xz}}' \hat{\mathbf{W}} \mathbf{S}_{\mathbf{xz}})^{-1}.
\]
\emph{Then:}
\begin{enumerate}[label=(\alph*)]
\item \emph{Under the null hypothesis $\mathrm{H}_0\colon \delta_\ell = \bar{\delta}_\ell$,}
  \[
    t_\ell
    = \frac{\sqrt{n}\,(\hat{\delta}_\ell(\hat{\mathbf{W}}) - \bar{\delta}_\ell)}
           {\sqrt{\bigl(\widehat{\avar(\hat{\boldsymbol{\delta}}(\hat{\mathbf{W}}))}\bigr)_{\!\ell\ell}}}
    = \frac{\hat{\delta}_\ell(\hat{\mathbf{W}}) - \bar{\delta}_\ell}{SE_\ell^*}
    \dto N(0,1),
  \]
  \emph{where
  $\bigl(\widehat{\avar(\hat{\boldsymbol{\delta}}(\hat{\mathbf{W}}))}\bigr)_{\!\ell\ell}$
  is the $(\ell,\ell)$ element of
  $\widehat{\avar(\hat{\boldsymbol{\delta}}(\hat{\mathbf{W}}))}$
  [which is $\avar(\hat{\delta}_\ell(\hat{\mathbf{W}}))$] and}
  \begin{equation}
    SE_\ell^* \text{ (robust standard error)}
    \equiv \sqrt{\frac{1}{n} \cdot
      \bigl(\widehat{\avar(\hat{\boldsymbol{\delta}}(\hat{\mathbf{W}}))}\bigr)_{\!\ell\ell}}.
    \label{eq:3.5.7}
  \end{equation}

\item \emph{Under the null hypothesis $\mathrm{H}_0\colon
  \mathbf{R}\boldsymbol{\delta} = \mathbf{r}$ where $\#\mathbf{r}$ is the number of restrictions
  (the dimension of $\mathbf{r}$) and $\mathbf{R}$ $(\#\mathbf{r} \times L)$ is of full row rank,}
  \begin{equation}
    W \equiv n \cdot (\mathbf{R}\hat{\boldsymbol{\delta}}(\hat{\mathbf{W}}) - \mathbf{r})'\,
    \bigl\{\mathbf{R}\bigl[\widehat{\avar(\hat{\boldsymbol{\delta}}(\hat{\mathbf{W}}))}\bigr]
    \mathbf{R}'\bigr\}^{-1}
    (\mathbf{R}\hat{\boldsymbol{\delta}}(\hat{\mathbf{W}}) - \mathbf{r})
    \dto \chi^2(\#\mathbf{r}).
    \label{eq:3.5.8}
  \end{equation}

\item \emph{Under the null hypothesis $\mathrm{H}_0\colon
  \mathbf{a}(\boldsymbol{\delta}) = \mathbf{0}$ such that $\mathbf{A}(\boldsymbol{\delta})$, the
  $\#\mathbf{a} \times L$ matrix of first derivatives of $\mathbf{a}$ (where $\#\mathbf{a}$ is the dimension
  of $\mathbf{a}$), is continuous and of full row rank,}
  \begin{equation}
    W \equiv n \cdot \mathbf{a}(\hat{\boldsymbol{\delta}}(\hat{\mathbf{W}}))'\,
    \bigl\{\mathbf{A}(\hat{\boldsymbol{\delta}}(\hat{\mathbf{W}}))\,
    \bigl[\widehat{\avar(\hat{\boldsymbol{\delta}}(\hat{\mathbf{W}}))}\bigr]\,
    \mathbf{A}(\hat{\boldsymbol{\delta}}(\hat{\mathbf{W}}))'\bigr\}^{-1}\,
    \mathbf{a}(\hat{\boldsymbol{\delta}}(\hat{\mathbf{W}}))
    \dto \chi^2(\#\mathbf{a}).
    \label{eq:3.5.9}
  \end{equation}
\end{enumerate}
\end{proposition}

For the Wald statistic for the nonlinear hypothesis, the same comment that we
made for Proposition~2.3 applies here: the numerical value of the Wald statistic is
not invariant to how the nonlinear restriction is represented by $\mathbf{a}(\cdot)$.

%% ─── Estimation of S ──────────────────────────────────────────────────────
\subsection*{Estimation of S}

We have already studied the estimation of
$\mathbf{S}$ $(= \E(\mathbf{g}_i \mathbf{g}_i') = \E(\varepsilon_i^2\, \mathbf{x}_i \mathbf{x}_i'))$
in Section~2.5.  The proposed estimator was, if adapted to the present estimation equation
$y_i = \mathbf{z}_i' \boldsymbol{\delta} + \varepsilon_i$,
\begin{equation}
  \hat{\mathbf{S}} \equiv \frac{1}{n} \sum_{i=1}^{n}
  \hat{\varepsilon}_i^2\, \mathbf{x}_i \mathbf{x}_i',
  \label{eq:3.5.10}
\end{equation}
where $\hat{\varepsilon}_i \equiv y_i - \mathbf{z}_i' \hat{\boldsymbol{\delta}}$ and
$\hat{\boldsymbol{\delta}}$ is consistent for $\boldsymbol{\delta}$.  The fourth-moment assumption
needed for this to be consistent for $\mathbf{S}$ is a generalization of Assumption~2.6.

\begin{assumption}[finite fourth moments]
\label{ass:3.6}
$\E[(x_{ik}\, z_{i\ell})^2]$ \emph{exists is finite for all}
$k$ $(= 1, \ldots, K)$ \emph{and} $\ell$ $(= 1, \ldots, L)$.
\end{assumption}

It is left as an analytical exercise to prove

\begin{proposition}[consistent estimation of S]
\label{prop:3.4}
\emph{Suppose the coefficient estimate $\hat{\boldsymbol{\delta}}$
used for calculating the residual $\hat{\varepsilon}_i$ for $\hat{\mathbf{S}}$ in \eqref{eq:3.5.10} is
consistent, and suppose $\mathbf{S} = \E(\mathbf{g}_i \mathbf{g}_i')$ exists and is finite.  Then,
under Assumptions~3.1, 3.2, and 3.6, $\hat{\mathbf{S}}$ given in \eqref{eq:3.5.10} is consistent
for $\mathbf{S}$.}
\end{proposition}

%% ─── Efficient GMM Estimator ──────────────────────────────────────────────
\subsection*{Efficient GMM Estimator}

Naturally, we wish to choose from the GMM estimators indexed by $\hat{\mathbf{W}}$ one that has
the least asymptotic variance.  The next proposition provides a choice of $\mathbf{W}$ that
minimizes the asymptotic variance.

\begin{proposition}[optimal choice of the weighting matrix]
\label{prop:3.5}
\emph{A lower bound for the asymptotic variance of the GMM estimators \eqref{eq:3.4.8}
indexed by $\hat{\mathbf{W}}$ is given by
$(\boldsymbol{\Sigma}_{\mathbf{xz}}' \mathbf{S}^{-1} \boldsymbol{\Sigma}_{\mathbf{xz}})^{-1}$.
The lower bound is achieved if $\hat{\mathbf{W}}$ is such that
$\mathbf{W}$ $(\equiv \plim\, \hat{\mathbf{W}}) = \mathbf{S}^{-1}$.\footnote{The condition that
$\mathbf{W} = \mathbf{S}^{-1}$ is sufficient but not necessary for efficiency.  A necessary and
sufficient condition is that there exists a matrix $\mathbf{C}$ such that
$\boldsymbol{\Sigma}_{\mathbf{xz}}' \mathbf{W} = \mathbf{C}\, \boldsymbol{\Sigma}_{\mathbf{xz}}' \mathbf{S}^{-1}$.
See Newey and McFadden (1994, p.~2165).}}
\end{proposition}

Because the asymptotic variance for any given $\hat{\mathbf{W}}$ is \eqref{eq:3.5.1}, this proposition
is proved if we can show that
\begin{equation}
  (\boldsymbol{\Sigma}_{\mathbf{xz}}' \mathbf{W} \boldsymbol{\Sigma}_{\mathbf{xz}})^{-1}\,
  \boldsymbol{\Sigma}_{\mathbf{xz}}' \mathbf{W} \mathbf{S} \mathbf{W}
  \boldsymbol{\Sigma}_{\mathbf{xz}}\,
  (\boldsymbol{\Sigma}_{\mathbf{xz}}' \mathbf{W} \boldsymbol{\Sigma}_{\mathbf{xz}})^{-1}
  \;\geq\;
  (\boldsymbol{\Sigma}_{\mathbf{xz}}' \mathbf{S}^{-1} \boldsymbol{\Sigma}_{\mathbf{xz}})^{-1}
  \label{eq:3.5.11}
\end{equation}
for any symmetric positive definite matrix $\mathbf{W}$.  Proving this algebraic result is left
as an analytical exercise.

A GMM estimator satisfying the efficiency condition that
$\plim\, \hat{\mathbf{W}} = \mathbf{S}^{-1}$ will be called the \textbf{efficient} (or \textbf{optimal})
\textbf{GMM estimator}.  Simply by replacing $\hat{\mathbf{W}}$ by $\hat{\mathbf{S}}^{-1}$ (which is
consistent for $\mathbf{S}^{-1}$) in the formulas of Proposition~3.1, we obtain
\begin{align}
  &\text{Efficient GMM estimator:} \quad
  \hat{\boldsymbol{\delta}}(\hat{\mathbf{S}}^{-1})
  = (\mathbf{S}_{\mathbf{xz}}' \hat{\mathbf{S}}^{-1} \mathbf{S}_{\mathbf{xz}})^{-1}\,
    \mathbf{S}_{\mathbf{xz}}' \hat{\mathbf{S}}^{-1} \mathbf{s}_{\mathbf{xy}},
  \label{eq:3.5.12} \\
  &\avar(\hat{\boldsymbol{\delta}}(\hat{\mathbf{S}}^{-1}))
  = (\boldsymbol{\Sigma}_{\mathbf{xz}}' \mathbf{S}^{-1}
     \boldsymbol{\Sigma}_{\mathbf{xz}})^{-1},
  \label{eq:3.5.13} \\
  &\widehat{\avar(\hat{\boldsymbol{\delta}}(\hat{\mathbf{S}}^{-1}))}
  = (\mathbf{S}_{\mathbf{xz}}' \hat{\mathbf{S}}^{-1}
     \mathbf{S}_{\mathbf{xz}})^{-1}.
  \label{eq:3.5.14}
\end{align}
With $\hat{\mathbf{W}} = \hat{\mathbf{S}}^{-1}$, the formulas for the robust $t$ and the Wald statistics
in Proposition~3.3 become
\begin{equation}
  t_\ell = \frac{\hat{\delta}_\ell(\hat{\mathbf{S}}^{-1}) - \bar{\delta}_\ell}
                 {SE_\ell^*},
  \label{eq:3.5.15}
\end{equation}
where $SE_\ell^*$ is the robust standard error given by
\[
  SE_\ell^* = \sqrt{\frac{1}{n} \cdot
    \bigl((\mathbf{S}_{\mathbf{xz}}' \hat{\mathbf{S}}^{-1}
    \mathbf{S}_{\mathbf{xz}})^{-1}\bigr)_{\!\ell\ell}},
\]
and
\begin{equation}
  W = n \cdot \mathbf{a}(\hat{\boldsymbol{\delta}}(\hat{\mathbf{S}}^{-1}))'\,
  \bigl\{\mathbf{A}(\hat{\boldsymbol{\delta}}(\hat{\mathbf{S}}^{-1}))\,
  (\mathbf{S}_{\mathbf{xz}}' \hat{\mathbf{S}}^{-1} \mathbf{S}_{\mathbf{xz}})^{-1}\,
  \mathbf{A}(\hat{\boldsymbol{\delta}}(\hat{\mathbf{S}}^{-1}))'\bigr\}^{-1}\,
  \mathbf{a}(\hat{\boldsymbol{\delta}}(\hat{\mathbf{S}}^{-1})).
  \label{eq:3.5.16}
\end{equation}

To calculate the efficient GMM estimator, we need the consistent estimator
$\hat{\mathbf{S}}$.  But Proposition~3.4 assures us that the $\hat{\mathbf{S}}$ based on any consistent
estimator of $\boldsymbol{\delta}$ is consistent for $\mathbf{S}$.  This leads us to the following
\textbf{two-step efficient GMM procedure}:

\medskip
\noindent\emph{Step~1:} Choose a matrix $\hat{\mathbf{W}}$ that converges in probability to a symmetric and positive
definite matrix, and minimize $J(\tilde{\boldsymbol{\delta}}, \hat{\mathbf{W}})$ over
$\tilde{\boldsymbol{\delta}}$ to obtain $\hat{\boldsymbol{\delta}}(\hat{\mathbf{W}})$.  There
is no shortage of such matrices $\hat{\mathbf{W}}$ (e.g., $\hat{\mathbf{W}} = \mathbf{I}$), but
usually we set $\hat{\mathbf{W}} = \mathbf{S}_{\mathbf{xx}}^{-1}$.  The resulting estimator
$\hat{\boldsymbol{\delta}}(\mathbf{S}_{\mathbf{xx}}^{-1})$ is the celebrated two-stage least
squares (as we will see in Section~3.8).  Use this to calculate the residual
$\hat{\varepsilon}_i \equiv y_i - \mathbf{z}_i' \hat{\boldsymbol{\delta}}(\hat{\mathbf{W}})$ and obtain
a consistent estimator $\hat{\mathbf{S}}$ of $\mathbf{S}$ by \eqref{eq:3.5.10}.

\medskip
\noindent\emph{Step~2:} Minimize
$J(\tilde{\boldsymbol{\delta}}, \hat{\mathbf{S}}^{-1})$ over $\tilde{\boldsymbol{\delta}}$.
The minimizer is the efficient GMM estimator.

%% ─── Asymptotic Power ──────────────────────────────────────────────────────
\subsection*{Asymptotic Power}

Like the coefficient estimator, the $t$ and Wald statistics depend on the choice of
$\mathbf{W}$.  It seems intuitively obvious that those statistics associated with the efficient GMM
estimator should be preferred in large samples.  This intuition can be formalized in
terms of asymptotic power introduced in Section~2.4.  Take, for example, the $t$-ratio
for testing $\mathrm{H}_0\colon \delta_\ell = \bar{\delta}_\ell$.  The $t$-ratio is written (reproducing
\eqref{eq:3.5.7}) as
\begin{equation}
  t_\ell = \frac{\sqrt{n}\,(\hat{\delta}_\ell(\hat{\mathbf{W}})
    - \bar{\delta}_\ell)}
    {\sqrt{\bigl(\widehat{\avar(\hat{\boldsymbol{\delta}}(\hat{\mathbf{W}}))}\bigr)_{\!\ell\ell}}}.
  \label{eq:3.5.17}
\end{equation}
The denominator converges to some finite number even when the null is false.
In contrast, the numerator explodes (its plim is infinite) when the null is false.
So the power under any fixed alternative approaches unity.  That is, the test is
consistent.  This is true for any choice of $\mathbf{W}$, so test consistency cannot be a basis
for choosing $\mathbf{W}$.

Next, consider a sequence of local alternatives subject to Pitman drift:
\begin{equation}
  \delta_\ell^{(n)} = \bar{\delta}_\ell + \frac{\gamma}{\sqrt{n}}
  \label{eq:3.5.18}
\end{equation}
for some given $\gamma \neq 0$.  Substituting \eqref{eq:3.5.18} into \eqref{eq:3.5.17}, the
$t$-ratio above can be rewritten as
\begin{equation}
  t_\ell = \frac{\sqrt{n}\,(\hat{\delta}_\ell(\hat{\mathbf{W}})
    - \delta_\ell^{(n)})}
    {\sqrt{\bigl(\widehat{\avar(\hat{\boldsymbol{\delta}}(\hat{\mathbf{W}}))}\bigr)_{\!\ell\ell}}}
  + \frac{\gamma}
    {\sqrt{\bigl(\widehat{\avar(\hat{\boldsymbol{\delta}}(\hat{\mathbf{W}}))}\bigr)_{\!\ell\ell}}}.
  \label{eq:3.5.19}
\end{equation}
Applying the same argument for deriving the asymptotic distribution of the $t$-ratio
in (2.4.4), we can show that $t_\ell \dto N(\mu, 1)$, where
\begin{equation}
  \mu \equiv \frac{\gamma}
    {\sqrt{\bigl(\avar(\hat{\boldsymbol{\delta}}(\hat{\mathbf{W}}))\bigr)_{\!\ell\ell}}}.
  \label{eq:3.5.20}
\end{equation}
If the significance level is $\alpha$, the asymptotic power is given by
$\Pr(|x| > t_{\alpha/2})$, where $x \sim N(\mu, 1)$ and $t_{\alpha/2}$ is the level-$\alpha$ critical
value.  Evidently, the larger is $|\mu|$, the higher is the asymptotic power.  And $|\mu|$
decreases with the asymptotic variance.  Therefore, the asymptotic power against local
alternatives is maximized by the efficient GMM estimator.

%% ─── Small-Sample Properties ──────────────────────────────────────────────
\subsection*{Small-Sample Properties}

Do these desirable asymptotic properties of the efficient GMM estimator and associated
test statistics carry over to their finite-sample distributions?  The efficient
GMM estimator uses $\hat{\mathbf{S}}^{-1}$, a function of estimated fourth moments, for
$\hat{\mathbf{W}}$.  Generally, it takes a substantially larger sample size to estimate fourth moments
reliably than to estimate first and second moments.  So we would expect the efficient GMM
estimator to have poorer small-sample properties than the GMM estimators that do
not use fourth moments for $\hat{\mathbf{W}}$.  The July 1996 issue of the \emph{Journal of Business
and Economic Statistics} has a number of papers examining the small-sample distribution
of GMM estimators and associated test statistics for various DGPs.  Their
overall conclusion is that the equally weighted GMM estimator with $\hat{\mathbf{W}} = \mathbf{I}$
generally outperforms the efficient GMM in terms of the bias and variance in finite
samples.  They also find that the size of the efficient Wald statistic in small samples
far exceeds the assumed significance level.  That is, if $\alpha$ is the assumed significance
level and $c_\alpha$ is the associated critical value so that
$\Pr(\chi^2 > c_\alpha) = \alpha$, the probability in finite samples that the Wald statistic is
greater than $c_\alpha$ far exceeds $\alpha$; the test rejects the null too often.  Unfortunately,
however, like most other small-sample studies, those studies fail to produce clear
quantitative guidance which the empirical researcher could follow.

%% ─── Questions for Review ──────────────────────────────────────────────────
\bigskip
\begin{center}\rule{0.6\textwidth}{0.4pt}\end{center}
\noindent\textsc{Questions for Review}
\medskip

\begin{enumerate}
\item Verify that all the results of Sections~2.3--2.5 are special cases of those of this
  section.  In particular, verify that \eqref{eq:3.5.1} reduces to (2.3.4).
  \textbf{Hint:} $\boldsymbol{\Sigma}_{\mathbf{xz}}$ is square if $\mathbf{z}_i = \mathbf{x}_i$.

\item (Singular $\mathbf{W}$) \; Suppose $\mathbf{W}$ is singular but
  $\boldsymbol{\Sigma}_{\mathbf{xz}}' \mathbf{W} \boldsymbol{\Sigma}_{\mathbf{xz}}$ is nonsingular.
  Verify that all the results of this section (except Proposition~3.5) remain valid.

\item (Asymptotically equivalent choice of $\hat{\mathbf{W}}$) \; Suppose
  $\hat{\mathbf{W}}_1 - \hat{\mathbf{W}}_2 \pto \mathbf{0}$.  Show that
  \[
    \sqrt{n}\,\hat{\boldsymbol{\delta}}(\hat{\mathbf{W}}_1)
    - \sqrt{n}\,\hat{\boldsymbol{\delta}}(\hat{\mathbf{W}}_2) \pto \mathbf{0}.
  \]
  \textbf{Hint:}
  \begin{align*}
    &\sqrt{n}\,\hat{\boldsymbol{\delta}}(\hat{\mathbf{W}}_1)
    - \sqrt{n}\,\hat{\boldsymbol{\delta}}(\hat{\mathbf{W}}_2) \\
    &\quad = \bigl[(\mathbf{S}_{\mathbf{xz}}' \hat{\mathbf{W}}_1 \mathbf{S}_{\mathbf{xz}})^{-1}
      \mathbf{S}_{\mathbf{xz}}' \hat{\mathbf{W}}_1
      - (\mathbf{S}_{\mathbf{xz}}' \hat{\mathbf{W}}_2 \mathbf{S}_{\mathbf{xz}})^{-1}
      \mathbf{S}_{\mathbf{xz}}' \hat{\mathbf{W}}_2\bigr]\, \sqrt{n}\, \bar{\mathbf{g}}.
  \end{align*}
  $\sqrt{n}\, \bar{\mathbf{g}}$ converges in distribution to a random variable.
  Use Lemma~2.4(b).

\item (Three-step GMM) \; Consider adding to the efficient two-step GMM procedure
  the following third step: Recompute $\hat{\mathbf{S}}$ by \eqref{eq:3.5.10}, but this time using
  the residual from the second step.  Calculate the GMM estimator with this
  recomputed $\hat{\mathbf{S}}$.  Is this estimator consistent?  Asymptotically normal?  Efficient?
  \textbf{Hint:} Verify by Proposition~3.4 that the recalculated $\hat{\mathbf{S}}$ is consistent for
  $\mathbf{S}$.  Does the asymptotic distribution of the GMM estimator depend on the choice of
  $\hat{\mathbf{W}}$ if $\plim_{n\to\infty} \hat{\mathbf{W}}$ is the same?

\item (When $\mathbf{z}_i$ is a strict subset of $\mathbf{x}_i$) \; Suppose $\mathbf{z}_i$
  is a strict subset of $\mathbf{x}_i$.  So $\mathbf{x}_i$ includes, in addition to the regressors
  (which are all predetermined), some other predetermined variables.  Are the efficient
  two-step GMM estimator and the OLS estimator numerically the same?
  [Answer: No.]

\item (Data matrix representation of efficient GMM) \; Let $\mathbf{B}$ be the $n \times n$
  diagonal matrix whose $i$-th element is $\hat{\varepsilon}_i^2$, where $\hat{\varepsilon}_i$ is the
  residual from the first-step consistent estimation.  That is,
  \[
    \mathbf{B} \equiv \begin{bmatrix} \hat{\varepsilon}_1^2 & & \\ & \ddots & \\
    & & \hat{\varepsilon}_n^2 \end{bmatrix}.
  \]
  Verify that
  \[
    \hat{\boldsymbol{\delta}}(\hat{\mathbf{S}}^{-1})
    = [\mathbf{Z}'\mathbf{X}(\mathbf{X}'\mathbf{B}\mathbf{X})^{-1}\mathbf{X}'\mathbf{Z}]^{-1}\,
      \mathbf{Z}'\mathbf{X}(\mathbf{X}'\mathbf{B}\mathbf{X})^{-1}\mathbf{X}'\mathbf{y},
  \]
  where $\mathbf{X}$, $\mathbf{y}$, and $\mathbf{Z}$ are data matrices for the instruments, the dependent
  variable, and the regressors (they are defined in Section~3.8 below).

\item (GLS interpretation of efficient GMM) \; Let $\mathbf{X}$, $\mathbf{Z}$, and $\mathbf{y}$ be as in
  the previous question.  Then the estimation equation can be written as
  $\mathbf{y} = \mathbf{Z}\boldsymbol{\delta} + \boldsymbol{\varepsilon}$.  Premultiply both sides by
  $\mathbf{X}'$ to obtain
  \[
    \mathbf{X}'\mathbf{y} = \mathbf{X}'\mathbf{Z}\boldsymbol{\delta}
    + \mathbf{X}'\boldsymbol{\varepsilon}.
  \]
  Taking $\mathbf{S}$ to be the variance matrix of $\mathbf{X}'\boldsymbol{\varepsilon}$ and
  $\hat{\mathbf{S}}$ to be its consistent estimate, apply FGLS.  Verify that the FGLS estimator is
  the efficient GMM estimator (the FGLS estimator was defined in Section~1.6).

\item (Linear combination of orthogonality conditions) \; Derive the efficient GMM
  estimator that exploits a linear combination of orthogonality conditions,
  \[
    \mathbf{A}\,\E(\mathbf{g}_i) = \mathbf{0},
  \]
\end{enumerate}
